% rejection-sampling.tex — \input'd from experiments.tex
%
% Rejection sampling pipeline for generating random bounded irredundant
% 4-polytopes. Contains a self-contained proposition + proof for the
% boundedness characterization (Proposition~\ref{prop:boundedness-iff}).
%
% Sources:
%   - Proposition and proof: standard convex geometry (recession cones,
%     positive span). Proof written by Claude, reviewed by Jörn.
%   - Acceptance rate data: real code execution (validation.rs).
%
% Structure:
%   - Sampling procedure (halfspace representation, 4-step validation)
%   - Proposition: bounded iff (span R^4) and (all triples pass)
%   - Proof: recession cone ↔ bounded, positive span ↔ (i)+(ii)
%   - Acceptance rate sweep (table + observations)

% Jörn: text approved (57e73d9) — from \subsection{Rejection sampling} to "Acceptance rate sweep"
\subsection{Rejection sampling acceptance rates}
\label{sec:rejection-sampling}

To probe Viterbo's conjecture computationally, we need large datasets of
random convex polytopes in~$\mathbb{R}^4$. We generate these by rejection
sampling: sample a candidate halfspace representation, then validate that it
defines a bounded, irredundant polytope.

\subsubsection*{Sampling procedure}

A polytope~$K \subset \mathbb{R}^4$ is specified in halfspace representation
\[
  K = \bigl\{\, x \in \mathbb{R}^4 \;\big|\; n_i \cdot x \le h_i,\;
  i = 1, \ldots, F \,\bigr\},
\]
where~$n_i \in S^3$ are outward unit normals and~$h_i > 0$ are heights
(positive heights ensure the origin lies in the interior).

For a given facet count~$F$ and height range~$[h_{\min}, h_{\max}]$, each
sample attempt proceeds as follows:
\begin{enumerate}
  \item Sample~$F$ independent unit normals uniformly on~$S^3$
    (via normalized standard Gaussian vectors in~$\mathbb{R}^4$).
  \item Sample~$F$ independent heights uniformly from~$[h_{\min}, h_{\max}]$.
  \item Validate the candidate~$(n_1, \ldots, n_F, h_1, \ldots, h_F)$:
  \begin{enumerate}
    \item \emph{No duplicate normals}: reject if~$\|n_i - n_j\| < \varepsilon$
      for any~$i \ne j$.
      % Same-direction only (n_i ≈ n_j). Antiparallel normals (n_i ≈ -n_j)
      % are valid opposite-facing facets, e.g. hypercube has ±e_i.
    \item \emph{Boundedness}: the polytope must be bounded
      (Proposition~\ref{prop:boundedness-iff} below). We first check that the
      normals span~$\mathbb{R}^4$ (via SVD rank), then enumerate all
      triples~$(i,j,k)$: the three normals share a 1D kernel direction~$d$
      (computed via the generalized cross product in~$\mathbb{R}^4$), and we
      verify that some normal has positive inner product with~$d$ and some
      normal has negative inner product with~$d$.
      % The degenerate case (three coplanar normals, kernel ≥ 2D) is skipped in
      % the code — probability zero for continuous distributions on S^3.
    \item \emph{Vertex enumeration}: a vertex lies at the intersection of
      exactly~$4$ facet hyperplanes. Solve all~$\binom{F}{4}$ linear systems,
      discard infeasible or duplicate solutions.
    \item \emph{Irredundancy}: each facet must be a genuine 3-dimensional
      face. Check that the vertices lying on each facet hyperplane span a
      3-dimensional affine subspace (via SVD).
  \end{enumerate}
  \item If all checks pass, accept the polytope.
\end{enumerate}

\begin{proposition}[Boundedness check]\label{prop:boundedness-iff}
Write $K = \{x : Ax \le h\}$ where~$A$ has rows~$n_1, \ldots, n_F$ and
$h_\ell > 0$ for all~$\ell$ (so $0 \in \operatorname{int} K$). Then $K$ is
bounded if and only if
\textup{(i)}~the normals span~$\mathbb{R}^4$, and
\textup{(ii)}~for every triple~$(i,j,k)$ whose normals have a 1-dimensional
kernel direction~$d$, some normal has positive inner product with~$d$ and some
has negative inner product with~$d$.
% When rank(A) <= 2, no triple has a 1D kernel, so the triple check is
% vacuously true but K is unbounded. The "span R^4" clause handles this.
\end{proposition}

\begin{proof}
The recession cone $\operatorname{rec}(K) = \{d : Ad \le 0\}$ satisfies
$\operatorname{rec}(K) \subset K$ (since $0 \in K$, any $d \in
\operatorname{rec}(K)$ gives a ray $\{td : t \ge 0\} \subset K$).
Conversely, if~$K$ is unbounded it contains a ray $\{x_0 + td : t \ge 0\}$
for some nonzero~$d$; substituting into $A(x_0 + td) \le h$ and letting
$t \to \infty$ forces $Ad \le 0$, so $d \in \operatorname{rec}(K)$.  Hence
\[
  K \text{ bounded}
  \;\Longleftrightarrow\;
  \operatorname{rec}(K) = \{0\}.
\]
Now $\operatorname{rec}(K) = \{0\}$ means there is no nonzero~$d$ with
$n_\ell \cdot d \le 0$ for all~$\ell$; equivalently, for every nonzero~$d$,
some $n_\ell \cdot d > 0$.  This is exactly the condition that the normals
\emph{positively span}~$\mathbb{R}^4$.  It remains to show that positive span
is equivalent to the conjunction of~\textup{(i)} and~\textup{(ii)}.

\medskip\noindent\emph{($\Rightarrow$) Positive span implies both conditions.}
Positive span implies span~$\mathbb{R}^4$ (a fortiori). For the triple check:
let~$d$ be a kernel direction of triple~$(i,j,k)$, so $n_i \cdot d = n_j
\cdot d = n_k \cdot d = 0$. By positive span, some~$n_\ell$ has $n_\ell \cdot
d > 0$; since the triple's own normals have zero dot product with~$d$, this
$n_\ell$ lies outside the triple. Applying the same argument to~$-d$ gives a
normal with negative inner product.
% QC: The kernel is 1-dimensional by hypothesis, so d is determined up to
% sign. The condition "some positive, some negative" is sign-symmetric,
% hence independent of the choice of generator.

\medskip\noindent\emph{($\Leftarrow$) Span~$\mathbb{R}^4$ and all triples
pass implies positive span.}
We prove the contrapositive: if positive span fails but the normals
span~$\mathbb{R}^4$, some triple fails.

Suppose $\operatorname{rec}(K) \neq \{0\}$, i.e., there exists a nonzero~$d$
with $n_\ell \cdot d \le 0$ for all~$\ell$. We produce a triple~$(i,j,k)$
whose kernel direction satisfies $n_\ell \cdot d \le 0$ for all~$\ell$
(so no normal blocks it).

Among all unit vectors in~$\operatorname{rec}(K)$, pick~$d$ maximizing the
number of orthogonal normals $S_0 = \{\ell : n_\ell \cdot d = 0\}$.  We claim
that~$\{n_\ell : \ell \in S_0\}$ contains three linearly independent normals.
If not, $\{n_\ell : \ell \in S_0\}$ spans a subspace of dimension~$\le 2$,
so its orthogonal complement~$W$ has $\dim W \ge 4 - 2 = 2$.  Since
$n_\ell \cdot d = 0$ for every $\ell \in S_0$, we have $d \in W$.  Because
$\dim W \ge 2$ and $d \ne 0$, there exists $v \in W$ not parallel to~$d$.
Since $v \perp n_\ell$ for all $\ell \in S_0$, we have
\[
  n_\ell \cdot (d + tv) = n_\ell \cdot d + t\,(n_\ell \cdot v)
  = \begin{cases}
    0 + t \cdot 0 = 0 & \text{if } \ell \in S_0, \\
    \underbrace{n_\ell \cdot d}_{<\,0} + t\,(n_\ell \cdot v) < 0
    & \text{if } \ell \notin S_0, \text{ for small } |t|.
  \end{cases}
\]
% QC: ℓ ∉ S₀ is non-empty because normals span R^4 ⟹ rank(A) = 4
% ⟹ ker(A) = {0}, so the nonzero vector d cannot satisfy n_ℓ · d = 0 for all ℓ.
(The case $\ell \notin S_0$ is non-empty: since the normals span~$\mathbb{R}^4$,
the matrix~$A$ has trivial kernel, so~$d \ne 0$ cannot be orthogonal to all
normals.)
So $d + tv \in \operatorname{rec}(K)$ for small~$|t|$.

Consider increasing~$|t|$ from zero.  The set
$T = \{t \in \mathbb{R} : n_\ell \cdot (d + tv) \le 0 \;\forall \ell\}$ is
a closed interval containing~$0$ (the intersection of half-lines
$\{t : n_\ell \cdot d + t\,n_\ell \cdot v \le 0\}$).  This interval is
bounded: if $d + tv \in \operatorname{rec}(K)$ for all~$t$, then
$n_\ell \cdot v = 0$ for every~$\ell$ with $n_\ell \cdot d < 0$; combined
with $n_\ell \cdot v = 0$ for $\ell \in S_0$, this gives $n_\ell \cdot v = 0$
for all~$\ell$, contradicting $\operatorname{rank}(A) = 4$.

So $T = [t^-, t^+]$ is a bounded closed interval, and at least one
endpoint~$t^*$ is nonzero.  At $t = t^*$, some new index
$\ell^* \notin S_0$ satisfies $n_{\ell^*} \cdot (d + t^* v) = 0$.  Set
$d' = d + t^* v$; this is nonzero because $d$ and~$v$ are not parallel.
Since $v \in W$, we still have $n_\ell \cdot d' = n_\ell \cdot d = 0$
for every $\ell \in S_0$.  Thus
$\{\ell : n_\ell \cdot d' = 0\} \supseteq S_0 \cup \{\ell^*\}$,
a strict superset of~$S_0$.  Since $d' \in \operatorname{rec}(K)$, the unit
vector $d' / \|d'\|$ contradicts the choice of~$d$ as maximizing~$|S_0|$.

So~$\{n_\ell : \ell \in S_0\}$ contains three linearly independent normals;
pick such a triple~$(i,j,k)$ from~$S_0$.  Since $n_i, n_j, n_k$ are linearly
independent in~$\mathbb{R}^4$, their common kernel $\ker(n_i, n_j, n_k)$ is
$(4 - 3)$-dimensional, so the triple falls under condition~\textup{(ii)}.
The direction~$d$ spans this kernel, and $d \in \operatorname{rec}(K)$ gives
$n_\ell \cdot d \le 0$ for all~$\ell$.  In particular, no normal has positive
inner product with~$d$, so triple~$(i,j,k)$ fails condition~\textup{(ii)}.
\end{proof}

\subsubsection*{Acceptance rate sweep}

We measured acceptance rates across a grid of parameters: facet
counts~$F \in \{5, 6, \ldots, 10\}$ and three height
ranges~$[h_{\min}, h_{\max}]$, from narrow ($[0.8, 1.2]$, heights nearly
equal) to wide ($[0.1, 5.0]$, heights spanning a 50$\times$ range). For each
of the $6 \times 3 = 18$ configurations, 1000~independent attempts were made
with a deterministic random seed.

\begin{table}[ht]
\centering
\begin{tabular}{r ccc}
\toprule
& \multicolumn{3}{c}{Acceptance rate} \\
\cmidrule(lr){2-4}
$F$ & $h \in [0.5, 2.0]$ & $h \in [0.1, 5.0]$ & $h \in [0.8, 1.2]$ \\
\midrule
 5 & 0.057 & 0.057 & 0.057 \\
 6 & 0.168 & 0.162 & 0.170 \\
 7 & 0.334 & 0.305 & 0.344 \\
 8 & 0.450 & 0.371 & 0.506 \\
 9 & 0.489 & 0.332 & 0.602 \\
10 & 0.504 & 0.283 & 0.729 \\
\bottomrule
\end{tabular}
\caption{Acceptance rates for rejection sampling of random 4-polytopes
  with~$F$ facets. Each cell reports the fraction of 1000 attempts that passed
  all validation checks.}
\label{tab:acceptance-rates}
\end{table}

\subsubsection*{Observations}

\paragraph{Acceptance increases with~$F$ for narrow height ranges.}
For~$h \in [0.8, 1.2]$, the rate climbs from~5.7\% at~$F = 5$ to~72.9\%
at~$F = 10$. More facets provide more directions to ``close off'' the
polytope, making boundedness easier to satisfy.

\paragraph{Acceptance is insensitive to height range at low~$F$.}
At~$F = 5$, all three height ranges give exactly the same rate (5.7\%).
This is not a coincidence: with only~5 facets, most random normal
configurations fail the boundedness check regardless of the heights chosen.
Since boundedness depends only on the normal directions, varying the height
range has no effect.
% The same seed produces the same normals for all three h-ranges of a given F.
% Identical acceptance at F=5 confirms that rejection is dominated by the
% boundedness check (which is height-independent), not irredundancy.

\paragraph{Wide height ranges hurt at high~$F$.}
At~$F = 10$, the rate drops from 72.9\% (for $h \in [0.8, 1.2]$) to 28.3\%
(for $h \in [0.1, 5.0]$). With~10 facets, boundedness is almost always
satisfied, so the irredundancy check becomes the bottleneck. When heights vary
over a wide range, some facets end up with too few incident vertices to span
a 3-dimensional face, causing rejection.

\paragraph{Practical consequence.}
For building datasets with~$F \ge 7$ and moderate height ranges, rejection
sampling is efficient (acceptance~$\ge 30\%$). Even the worst case ($F = 5$,
5.7\% acceptance) requires only~${\sim}18$ attempts per accepted polytope,
which is practical for datasets of a few thousand samples.
